NLP healthcare applications have almost exclusively been clinician facing, even though patients have had increasing access, by law and by practice, to their own clinical texts. A systematic review of NLP on patient-authored text found that even work using patient-generated data has been oriented primarily toward clinical surveillance and institutional quality improvement; the patient who generated the text rarely benefits directly from the analysis. Yet the tools themselves are not inherently clinician-facing or applicable only to clinical text. This talk explores how NLP can serve as a mirror through which patients see themselves reflected in their own data. Drawing on a fifty-year longitudinal auto-ethnographic corpus blending memoir, personal journals, medical records, wearable data, and other artifacts, I juxtapose what clinical records document with what personal narratives reveal, illustrating meaning making across diverse health circumstances using tools ranging from psycholinguistic text analysis to generative AI. I close by suggesting that NLP-enabled reflections can meaningfully support patient meaning making, and by sketching open questions for patient-facing NLP research